% CLASSE DO DOCUMENTO
% ======================================================================================================================
% Define a classe abntex2, baseada nas normas de formatação
% de trabalhos acadêmicos da ABNT.
%
% Opções:
%   12pt     : tamanho da fonte principal do documento;
%   article  : configura o documento no formato de artigo;
%   oneside  : imprime o documento apenas em um lado da folha;
%   a4paper  : define o tamanho do papel como A4;
%   english  : carrega suporte ao idioma inglês;
%   brazil   : define o português brasileiro como idioma principal, conforme a configuração da classe.
\documentclass[
    12pt,
    article,
    oneside,
    a4paper,
    english,
    brazil
]{abntex2}


% PACOTES DE TIPOGRAFIA E FORMATAÇÃO
% ======================================================================================================================
% Utiliza a família de fontes Latin Modern, uma versão aprimorada das fontes Computer Modern tradicionais do LaTeX.
\usepackage{lmodern}
% Melhora o espaçamento entre caracteres e palavras, contribuindo para uma composição tipográfica mais uniforme.
\usepackage{microtype}
% Aplica recuo de parágrafo também ao primeiro parágrafo após títulos de seções e outros elementos.
\usepackage{indentfirst}

% PACOTES PARA MATEMÁTICA
% ======================================================================================================================
% Disponibiliza ambientes e comandos matemáticos avançados, incluindo equações, alinhamentos e expressões matemáticas.
\usepackage{amsmath}
% Disponibiliza fontes para conjuntos e símbolos matemáticos em negrito, caligráficos e outros estilos.
\usepackage{amsfonts}
% Acrescenta símbolos matemáticos adicionais, como símbolos de relações, operadores e conjuntos numéricos.
\usepackage{amssymb}
% Complementa o amsmath com recursos matemáticos adicionais, como comandos para delimitação e construção de expressões.
\usepackage{mathtools}
% Permite definir ambientes para axiomas, definições, teoremas, lemas, proposições, corolários e observações.
\usepackage{amsthm}

% CITAÇÕES E REFERÊNCIAS BIBLIOGRÁFICAS
% ======================================================================================================================
% Configura o sistema de citações da classe abntex2. A opção alf utiliza o sistema autor-data, de acordo com o estilo
% bibliográfico correspondente da ABNT.
\usepackage[alf]{abntex2cite}

% HIPERLINKS E METADADOS DO ARQUIVO PDF
% ======================================================================================================================
% O pacote hyperref cria links clicáveis no PDF para referências cruzadas, citações, URLs e outros elementos.
%
% O abntex2 pode carregar esse pacote internamente. Por isso, antes de carregá-lo explicitamente, verifique se ele já
% foi carregado pela classe ou por outro pacote. Caso já esteja disponível, não é necessário carregá-lo novamente.
\usepackage{hyperref}

% Configura a aparência dos links e os metadados do PDF.
\hypersetup{
    % Ativa a exibição dos links com cores, sem molduras.
    colorlinks = true,
    % Mantém pretos os links internos, como os de seções e referências cruzadas.
    linkcolor  = black,
    % Mantém pretos os links associados às citações.
    citecolor  = black,
    % Exibe URLs em azul, facilitando sua identificação.
    urlcolor   = blue,
    % Define o título exibido nas propriedades do PDF.
    pdftitle   = {Fundamentos da Teoria Axiomática dos Conjuntos},
    % Identifica a autoria do documento nos metadados.
    pdfauthor  = {Danilo Marto de Carvalho},
    % Informa o assunto principal do documento.
    pdfsubject = {Teoria dos Conjuntos},
    % Registra palavras-chave para identificação do conteúdo.
    pdfkeywords = {axiomas, ZF, teoria dos conjuntos, lógica matemática}
}

% AMBIENTES PARA AXIOMAS E DEFINIÇÕES
% ======================================================================================================================
% O estilo definition utiliza, em geral, títulos em negrito e corpo do texto em fonte romana.
\theoremstyle{definition}
% Cria o ambiente axioma, numerado por seção. Exemplo: Axioma 1.1, Axioma 1.2, Axioma 2.1 etc.
\newtheorem{axioma}{Axioma}[section]
% Cria o ambiente definicao, também numerado por seção.
\newtheorem{definicao}{Definição}[section]

% AMBIENTES PARA RESULTADOS MATEMÁTICOS
% ======================================================================================================================
% O estilo plain utiliza, em geral, título em negrito e corpo do texto em itálico, conforme a convenção tradicional de
% apresentação de resultados matemáticos.
\theoremstyle{plain}
% Cria o ambiente teorema, cuja numeração é reiniciada a cada seção.
\newtheorem{teorema}{Teorema}[section]
% O lema compartilha o contador dos teoremas. Assim, teoremas, lemas, corolários e proposições formam uma sequência
% única de numeração dentro de cada seção.
\newtheorem{lema}[teorema]{Lema}
% O corolário compartilha o contador dos teoremas.
\newtheorem{corolario}[teorema]{Corolário}
% A proposição também compartilha o contador dos teoremas.
\newtheorem{proposicao}[teorema]{Proposição}

% AMBIENTE PARA EXEMPLOS
% ======================================================================================================================
% Retorna ao estilo definition para que os exemplos sejam apresentados com título destacado e corpo romano.
\theoremstyle{definition}
% Cria o ambiente exemplo, numerado por seção.
\newtheorem{exemplo}{Exemplo}[section]

% AMBIENTE PARA OBSERVAÇÕES
% ======================================================================================================================
% O estilo remark é apropriado para observações e comentários, geralmente apresentados com corpo em fonte romana.
\theoremstyle{remark}
% Cria o ambiente observacao, numerado por seção.
\newtheorem{observacao}{Observação}[section]

% ELEMENTOS DE IDENTIFICAÇÃO DO DOCUMENTO
% ======================================================================================================================
% Define o título principal do artigo.
\titulo{Fundamentos da Teoria dos Conjuntos}
% Define o título em outro idioma.
\tituloestrangeiro{Fundamentals of Set Theory}
% Define o autor e acrescenta uma nota de rodapé com o endereço de e-mail para contato.
\autor{
    Danilo Marto de Carvalho%
    \thanks{E-mail: \texttt{carvalho.dm@proton.me}}
}
% Informa o local associado ao documento.
\local{Brasil}
% Utiliza a data atual de compilação.
\date{\today}

% INÍCIO DO DOCUMENTO
% ======================================================================================================================

\begin{document}
% Gera o título a partir das informações declaradas anteriormente, como título, autor e data.
\maketitle

% INCLUSÃO DOS CAPÍTULOS OU SEÇÕES DO ARTIGO
% ======================================================================================================================

% Inclui o conteúdo do arquivo referente ao Axioma da Extensão
\section{O Axioma da Extensão e os Conceitos Primitivos}\label{sec:axioma-extensao}

Na abordagem puramente axiomática da Teoria dos Conjuntos (TC), os conceitos de \textbf{conjunto}, \textbf{pertinência} e \textbf{igualdade} desempenham papéis fundamentais na estruturação da teoria. O conceito de conjunto e a relação de pertinência são considerados primitivos, enquanto a igualdade é um símbolo lógico primitivo herdado da Lógica de Primeira Ordem (CQC), cujas propriedades são estabelecidas pelos princípios lógicos correspondentes. Adicionalmente, a TC é uma teoria pura, na qual todos os objetos do domínio são, por definição, conjuntos.

A relação de pertinência entre dois conjuntos $x$ e $y$ é denotada por $x \in y$ (\textit{``$x$ pertence a $y$''}) e dizemos que $x$ é \textbf{elemento de $y$}, enquanto sua negação é representada por $x \notin y$. Analogamente, a igualdade é representada por $x = y$ (\textit{``$x$ é igual a $y$''}), tendo como negação a expressão $x \neq y$.

Como o conceito de conjunto e a relação de pertinência são primitivos, eles não recebem definições formais explícitas na linguagem da teoria. Seu significado e comportamento são determinados indiretamente pelos axiomas que regem a relação entre os objetos do domínio. A igualdade, por sua vez, possui suas propriedades lógicas gerais estabelecidas pela lógica subjacente, mas sua relação com a pertinência é também condicionada pelos axiomas específicos da TC.

O primeiro passo para estabelecer um critério de identidade entre conjuntos é dado pelo Axioma da Extensão, segundo o qual dois conjuntos que possuem exatamente os mesmos elementos são necessariamente iguais:

\begin{axioma}[Extensão]\label{ax:extensao}
Se dois conjuntos possuem exatamente os mesmos elementos, então são iguais. Formalmente,
\begin{equation}\label{eq:A1}
    \forall x \forall y \bigl[
        \forall z (z \in x \leftrightarrow z \in y)
        \rightarrow x = y
    \bigr].
\end{equation}
\end{axioma}

Note que a Fórmula~\eqref{eq:A1} expressa apenas uma das direções da caracterização usual da igualdade entre conjuntos: se dois conjuntos possuem exatamente os mesmos elementos, então são iguais. A direção inversa, isto é, a afirmação de que conjuntos iguais possuem os mesmos elementos, decorre dos princípios lógicos da igualdade, em particular de sua propriedade de substitutividade. Com efeito, se $x = y$, então, para qualquer conjunto $z$, vale $z \in x \leftrightarrow z \in y$. Portanto, em conjunto com os princípios lógicos da igualdade, o Axioma da Extensão permite concluir que dois conjuntos são iguais se, e somente se, possuem exatamente os mesmos elementos.

A relevância desse axioma reside no fato de que sua afirmação não decorre apenas da lógica de primeira ordem com igualdade. Embora a lógica estabeleça as propriedades formais da igualdade, ela não exige, por si só, que dois objetos com os mesmos elementos sejam idênticos. Essa exigência é introduzida especificamente pelo Axioma da Extensão, que impõe uma condição fundamental à relação de pertinência: conjuntos que possuem exatamente os mesmos elementos devem ser identificados como um único conjunto. Como observa Halmos~\cite[p.~4]{halmos2001}, \textit{``[...] o axioma da extensão não é apenas uma propriedade logicamente necessária sobre a igualdade, mas uma propriedade não trivial sobre a pertinência. [...]''}, pois estabelece que a igualdade entre conjuntos é determinada exclusivamente pela coincidência de seus elementos.

Dessa forma, as propriedades clássicas da igualdade — \textbf{reflexividade} ($x = x$), \textbf{simetria} ($x = y \rightarrow y = x$) e \textbf{transitividade} ($x = y \land y = z \rightarrow x = z$), atuam como infraestrutura lógica preexistente. O Axioma da Extensão completa esse quadro ao conectar diretamente essa estrutura lógica à relação primitiva de pertinência.

% Inclui o conteúdo do arquivo referente ao Axioma da Separação.
\section{O Esquema Axiomático da Especificação}
\label{sec:axioma-separacao}

Antes de apresentarmos o Esquema Axiomático da Especificação, é conveniente formalizarmos a relação de inclusão entre conjuntos, uma vez que o axioma a seguir pressupõe essa noção e a utiliza de modo essencial em sua caracterização.

\begin{definicao}[Inclusão e Subconjunto]
\label{def:inclusao}
Um conjunto $x$ é chamado de \textbf{subconjunto} de um conjunto $y$ (ou dito estar contido em $y$) se, e somente se, todos os elementos de $x$ também forem elementos de $y$.
\begin{equation}\label{eq:inclusao}
    x \subseteq y \leftrightarrow \forall z \bigl(z \in x \rightarrow z \in y\bigr).
\end{equation}
\end{definicao}

Lemos a notação $x \subseteq y$ como \textit{``$x$ está contido em $y$''}, \textit{``$x$ é subconjunto de $y$''} ou \textit{``$x$ é parte de $y$''}. Quando $x$ não é subconjunto de $y$ — isto é, quando existe ao menos um elemento de $x$ que não pertence a $y$ —, escrevemos $x \nsubseteq y$. A relação expressa na Equação~\eqref{eq:inclusao} pode ser escrita, na direção oposta, como $y \supseteq x$, lida como \textit{``$y$ contém $x$''}; sua negação é representada por $y \nsupseteq x$.

Quando um conjunto $x$ é subconjunto de $y$, mas não é igual a $y$, dizemos que $x$ é um \textbf{subconjunto próprio} de $y$ (ou que $x$ está contido propriamente em $y$):
\begin{equation}\label{eq:subconjunto-proprio}
    x \subsetneq y \leftrightarrow (x \subseteq y \land x \neq y).
\end{equation}

Vale destacar que há uma variação de convenções na literatura matemática quanto ao símbolo de inclusão. Na notação moderna e predominante em Teoria dos Conjuntos (adotada neste trabalho), utiliza-se $\subseteq$ para a inclusão ampla, de modo que vale a disjunção
\begin{equation*}
    x \subseteq y \;\leftrightarrow\; (x \subsetneq y \lor x = y).
\end{equation*}
Em contrapartida, alguns autores clássicos (como Halmos) utilizam $\subset$ para indicar a inclusão ampla, reservando $\subsetneq$ para a inclusão estrita; em outros textos, $\subset$ designa justamente a inclusão estrita. Como o símbolo $\subset$ é ambíguo na literatura, evitamo-lo neste trabalho em favor das notações $\subseteq$ e $\subsetneq$.

\begin{teorema}[Propriedades da Inclusão]
\label{teo:propriedades-inclusao}
A relação de inclusão entre conjuntos define uma ordem parcial ampla, satisfazendo as seguintes propriedades:
\begin{itemize}
    \item \textit{Reflexiva:} $\forall x\, (x \subseteq x)$;
    \item \textit{Antissimétrica:} $\forall x\, \forall y\, [(x \subseteq y \land y \subseteq x) \rightarrow x = y]$;
    \item \textit{Transitiva:} $\forall x\, \forall y\, \forall z\, [(x \subseteq y \land y \subseteq z) \rightarrow x \subseteq z]$.
\end{itemize}
\end{teorema}

\begin{proof}
\noindent\textit{(Reflexividade)} Para todo conjunto $x$ e todo elemento $z$, a implicação $z \in x \rightarrow z \in x$ é uma tautologia da lógica proposicional. Logo, $\forall z\, (z \in x \rightarrow z \in x)$ é verdadeira, o que, pela Definição~\ref{def:inclusao}, estabelece $x \subseteq x$.

\medskip
\noindent\textit{(Antissimetria)} Assumimos $x \subseteq y$ e $y \subseteq x$. Pela Definição~\ref{def:inclusao}, temos
\begin{equation*}
    \forall z\, (z \in x \rightarrow z \in y)
    \;\land\;
    \forall z\, (z \in y \rightarrow z \in x).
\end{equation*}
Pela distribuição do quantificador universal sobre a conjunção, obtemos
\begin{equation*}
    \forall z\, \bigl[(z \in x \rightarrow z \in y)
        \land (z \in y \rightarrow z \in x)\bigr].
\end{equation*}
Pela equivalência tautológica $(P \rightarrow Q) \land (Q \rightarrow P) \equiv (P \leftrightarrow Q)$, derivamos $\forall z\, (z \in x \leftrightarrow z \in y)$. Pelo Axioma da Extensão, concluímos que $x = y$.

\medskip
\noindent\textit{(Transitividade)} Suponha $x \subseteq y$ e $y \subseteq z$. Pela Definição~\ref{def:inclusao},
\begin{equation*}
    \forall w\, (w \in x \rightarrow w \in y)
    \;\land\;
    \forall w\, (w \in y \rightarrow w \in z).
\end{equation*}
Unificando sob um único quantificador universal, temos
\begin{equation*}
    \forall w\, \bigl[(w \in x \rightarrow w \in y)
        \land (w \in y \rightarrow w \in z)\bigr].
\end{equation*}
Pelo silogismo hipotético da lógica proposicional, $[(P \rightarrow Q) \land (Q \rightarrow R)] \rightarrow (P \rightarrow R)$, deduzimos $\forall w\, (w \in x \rightarrow w \in z)$. Pela Definição~\ref{def:inclusao}, concluímos que $x \subseteq z$.
\end{proof}

Com a relação de inclusão devidamente formalizada, podemos enunciar o axioma propriamente dito. Recorde que, na lógica de primeira ordem, assume-se que o domínio de discurso é não vazio; assim, existe ao menos um conjunto no universo, o que legitima o uso do Esquema da Separação a seguir.

\begin{axioma}[Especificação]
\label{ax:separacao}
Para cada conjunto $x$ e para cada fórmula bem formada $\varphi(z)$ na linguagem da teoria, existe um conjunto $y$ cujos elementos são precisamente os elementos $z \in x$ que satisfazem a propriedade $\varphi(z)$. Formalmente,
\begin{equation}\label{eq:axioma-separacao}
    \forall x\, \exists y\, \forall z\,
        \bigl[\, z \in y \;\leftrightarrow\;
            (z \in x \land \varphi(z)) \,\bigr].
\end{equation}
\end{axioma}

Do ponto de vista da Lógica de Primeira Ordem, a Equação~\eqref{eq:axioma-separacao} não representa um único axioma, mas sim um \textbf{esquema axiomático}: uma coleção infinitamente numerável de axiomas, obtida ao instanciar a fórmula $\varphi$ por qualquer condição sintaticamente válida. Tendo como única restrição indispensável que a variável $y$ não ocorra livre em $\varphi$, impedindo definições autorreferenciais na própria sintaxe da linguagem. Além disso, pela Definição~\ref{def:inclusao}, a estrutura do Esquema da Separação garante inerentemente que o conjunto resultante $y$ é um subconjunto de $x$, isto é, $y \subseteq x$.

Assim como o Axioma da Extensão estabelece um critério de identidade entre conjuntos, o Esquema da Separação estabelece um critério de \emph{existência} de subconjuntos: ele permite extrair, de um conjunto já dado, um novo conjunto determinado por uma propriedade arbitrária. Note, contudo, que o esquema não autoriza a formação de um conjunto a partir de uma propriedade qualquer tomada isoladamente — ele exige sempre um conjunto de referência $x$ do qual os elementos são separados. Essa restrição é precisamente o que impede a derivação de paradoxos como o de Russell, como veremos adiante.

\begin{teorema}[Unicidade da Separação]
\label{teo:unicidade-separacao}
O conjunto $y$ obtido pelo Esquema Axiomático da Separação a partir de um conjunto $x$ e de uma propriedade $\varphi$ é único.
\end{teorema}

\begin{proof}
O Esquema da Separação assegura a existência de $y$. Para provar a unicidade, suponha que existam dois conjuntos, $y_1$ e $y_2$, satisfazendo a condição para um mesmo conjunto $x$ e uma mesma fórmula $\varphi$. Então, para todo $z$,
\begin{equation*}
    z \in y_1 \leftrightarrow (z \in x \land \varphi(z))
    \quad \text{e} \quad
    z \in y_2 \leftrightarrow (z \in x \land \varphi(z)).
\end{equation*}
Pela transitividade e simetria da equivalência lógica ($\leftrightarrow$), deduzimos $\forall z\, (z \in y_1 \leftrightarrow z \in y_2)$. Pelo Axioma da Extensão, concluímos que $y_1 = y_2$.
\end{proof}

Garantida a unicidade pelo Axioma da Extensão, torna-se legítimo introduzir a notação formal por compreensão restrita:
\begin{equation*}
    y = \{z \in x \mid \varphi(z)\}.
\end{equation*}

A seguir, apresentamos duas aplicações fundamentais do Esquema da Separação, que ilustram tanto o seu poder quanto o seu papel na prevenção de paradoxos.

\subsection{Exemplo 1: Existência e Unicidade do Conjunto Vazio}
\label{sec:conjunto-vazio}

Garantida a existência de ao menos um conjunto no universo de discurso, o Esquema da Separação permite provar a existência de um conjunto que não possui elementos.

\begin{teorema}[Existência e Unicidade do Conjunto Vazio]
\label{teo:conjunto-vazio}
Existe um único conjunto que não possui elementos.
\end{teorema}

\begin{proof}
Seja $x$ um conjunto qualquer, cuja existência é garantida pelo fato de o domínio de discurso da lógica de primeira ordem ser não vazio. Escolhemos como fórmula a contradição lógica $\varphi(z) \equiv z \neq z$. Aplicando o Esquema da Separação ao conjunto $x$ com a fórmula $\varphi$, obtemos a existência de um conjunto $y$ tal que
\begin{equation*}
    \forall z\, \bigl[z \in y \leftrightarrow (z \in x \land z \neq z)\bigr].
\end{equation*}
Como $z \neq z$ é uma falsidade tautológica para todo $z$, a conjunção $(z \in x \land z \neq z)$ é sempre falsa. Logo, a proposição $z \in y$ é falsa para qualquer $z$, o que implica $\forall z\, (z \notin y)$. A unicidade de $y$ decorre diretamente do Teorema~\ref{teo:unicidade-separacao}.
\end{proof}

Pelo Teorema~\ref{teo:conjunto-vazio}, o conjunto sem elementos é único. Devido à sua relevância fundamental na teoria, ele recebe um símbolo próprio: $\emptyset$ (ou simplesmente $\{\}$).

Um resultado imediato sobre o conjunto vazio é que ele está contido em qualquer conjunto do domínio.

\begin{teorema}[Vacuidade da Inclusão do Conjunto Vazio]
\label{teo:vacuidade}
Qualquer que seja o conjunto $x$, o conjunto vazio é subconjunto de $x$.
\begin{equation}\label{eq:vacuidade}
    \forall x\, (\emptyset \subseteq x).
\end{equation}
\end{teorema}

\begin{proof}
Suponha, por contradição, que exista algum conjunto $x$ tal que $\emptyset \nsubseteq x$. Pela negação da Definição~\ref{def:inclusao}, isto implica que deve existir ao menos um objeto $z$ tal que $z \in \emptyset$ e $z \notin x$. No entanto, pela definição do conjunto vazio, a proposição $z \in \emptyset$ é falsa para todo $z$, o que contradiz a existência desse elemento. Portanto, a negação é falsa, concluindo-se que $\emptyset \subseteq x$ para todo conjunto $x$.
\end{proof}

\subsection{Exemplo 2: Inexistência do Conjunto Universo}
\label{sec:inexistencia-conjunto-universo}

Uma das consequências mais marcantes do Esquema da Separação é que ele impede a existência de um conjunto que contenha todos os objetos do universo.

\begin{teorema}[Inexistência do Conjunto Universo]
\label{teo:inexistencia-conjunto-universo}
Não existe um conjunto universo, isto é, não existe um conjunto $V$ tal que todo conjunto pertença a $V$.
\end{teorema}

\begin{proof}
Suponha, por contradição, que exista um conjunto universo $V$, ou seja, um conjunto tal que $\forall z\, (z \in V)$.

Aplicamos o Esquema da Separação ao conjunto $V$ utilizando a fórmula $\varphi(z) \equiv z \notin z$. Isto nos dá a existência de um conjunto $y$ definido por
\begin{equation*}
    y = \{z \in V \mid z \notin z\}.
\end{equation*}
Por construção, para todo $z$, temos $z \in y \leftrightarrow (z \in V \land z \notin z)$. Como assumimos que $V$ é o conjunto universo, a condição $z \in V$ é verdadeira para todo $z$; pela tautologia $(\text{Verdadeiro} \land P) \equiv P$, a equivalência simplifica-se para
\begin{equation*}
    \forall z\, (z \in y \leftrightarrow z \notin z).
\end{equation*}
Como $y$ é ele próprio um conjunto e $V$ contém todos os conjuntos, $y$ está no domínio de quantificação da variável $z$. Podemos, portanto, instanciar $z := y$, obtendo
\begin{equation*}
    y \in y \leftrightarrow y \notin y,
\end{equation*}
o que é uma contradição lógica da forma $P \leftrightarrow \neg P$. Pelo método de redução ao absurdo, concluímos que a premissa de que $V$ existe é falsa. Portanto, não existe o conjunto de todos os conjuntos.
\end{proof}


% REFERÊNCIAS BIBLIOGRÁFICAS
% ======================================================================================================================

% Define o estilo bibliográfico autor-data compatível com o pacote abntex2cite.
\bibliographystyle{abntex2-alf}
% Indica o arquivo que contém as entradas bibliográficas.
\bibliography{referencias}

% FIM DO DOCUMENTO
% ======================================================================================================================
\end{document}
